\section*{Chapitre \ref{ch:b2:palaeoenvironments} --- Fossiles et paléoenvironnements}

\begin{solution}{exo:b2:palaeoenvironments:1}
(a) $112/135 = 0.830$ et $112/150 = 0.747$. (b) Les coquilles mortes du fond s'accumulent pendant des années, voire
davantage~: elles comprennent des espèces qui y ont vécu plus tôt ou, plus rarement, qui y ont été apportées,
si bien que la \omterm{def:b2:palaeoenvironments:assemblages}{thanatocénose} fait la moyenne de nombreuses générations de la communauté. (c) Les groupes à corps
mou, surtout détritivores et carnivores comme les vers~: sans squelette, ils se décomposent avant l'enfouissement.
\end{solution}

\begin{solution}{exo:b2:palaeoenvironments:2}
(a) $5 \times 165 = \qty{825}{ka}$, soit environ \qty{0.8}{Ma}, dans les 0.7
à 1~Ma. (b) $183.73 - 182.77 = \qty{0.96}{Ma}$, et $0.96/0.165 = 5.8$~: soit
environ six horizons. (c) Une lignée change en beaucoup moins de 1~Ma~; seuls
les horizons séparent les formes successives, tandis qu'une zone les regroupe.
\end{solution}

\begin{solution}{exo:b2:palaeoenvironments:3}
(a) \omterm{def:b2:palaeoenvironments:ichnofacies}{Ichnofaciès} de Nereites~: le fond de mer profonde d'un système de turbidites, au-delà de la rupture
de pente, où la nourriture est rare. (b) \omterm{def:b2:palaeoenvironments:ichnofacies}{Ichnofaciès} de Skolithos~: estran ou avant-côte. (c) \omterm{def:b2:palaeoenvironments:ichnofacies}{Ichnofaciès}
de Cruziana~: au large, entre la base de la vague de beau temps et la base de la vague de tempête. (d)
\omterm{def:b2:palaeoenvironments:ichnofacies}{Ichnofaciès} de Zoophycos~: au Mésozoïque, plateau inférieur ou talus, sous la base de la vague de tempête.
\end{solution}

\begin{solution}{exo:b2:palaeoenvironments:4}
(a) Le lit A renferme des groupes sténohalins~: eau de mer normale. Le lit B, pauvre en espèces
mais riche en individus, avec des débris végétaux et aucun des groupes sténohalins, suggère une
eau saumâtre, par exemple un estuaire. (b) Les groupes manquants sont plusieurs, fréquents et
bien conservés dans des couches marines du même âge~: leur absence conjointe est un motif, alors
qu'un fossile peut manquer par hasard ou par mauvaise conservation. (c) Le sédiment (chenaux
de marée ou fluviaux, enduits de boue, racines), les \omterm{def:b2:palaeoenvironments:fossil-kinds}{traces fossiles} (un \omterm{def:b2:palaeoenvironments:ichnofacies}{ichnofaciès} à Cruziana
d'espèces peu nombreuses), et le passage latéral du lit B vers des dépôts marins ou fluviaux.
\end{solution}

\begin{solution}{exo:b2:palaeoenvironments:5}
(a) $P = 1/2$ lorsque $kh/s = \ln 2$~: $s = 0.05 \times 0.1/0.693 = \qty{0.0072}{m.an^{-1}}$, soit \qty{7.2}{mm}
par an. (b) Diviser $s$ par deux double l'exposant : $P = (1/2)^2 = 0.25$, soit la moitié de la valeur précédente.
(c) $k$ est divisé par deux : $P = e^{-2.5} = 0.082$ au lieu de $e^{-5} = 0.0067$, soit $e^{2.5} = 12$ fois plus de
chances. (d) Les coquilles épaisses et robustes sont surreprésentées, d'autant plus que l'enfouissement est lent.
\end{solution}

\begin{solution}{exo:b2:palaeoenvironments:6}
(a) La borne supérieure vaut $183.73 + 0.35 = \qty{184.08}{Ma}$~: la valeur de la charte, \qty{184.2}{Ma}, se situe
en dehors, d'un écart de \qty{0.12}{Ma}. (b) $183.73 - 182.77 = \qty{0.96}{Ma}$. (c) $6/0.96 =
\qty{6.25}{m}$ par million d'années. (d) Une biozone est définie par des fossiles, dont les apparitions et les
disparitions sont presque simultanées sur de vastes régions~; l'épaisseur de roche déposée pendant ce temps
dépend du taux local d'accumulation.
\end{solution}

\begin{solution}{exo:b2:palaeoenvironments:7}
(a) $0.49\,e^{-1046/3704} = 0.369$ et $0.49\,e^{-1166/3704} = 0.358$~; moyenne $0.364$. (b) $120
\times (1 - 0.364) = \qty{76.3}{m}$. (c) $76.3/(1 - 0.49) = \qty{149.6}{m}$. (d) Le facteur
suppose que toute la couche restaurée a la porosité $\varphi_0$ de la surface~; or une couche de
\qty{150}{m} d'épaisseur en surface est moins poreuse à sa base qu'à son sommet, si bien que les mêmes grains
remplissent moins que \qty{149.6}{m}.
\end{solution}

\begin{solution}{exo:b2:palaeoenvironments:8}
(a) Un volume unitaire contient $\varphi$ d'eau et $1 - \varphi$ de grains~: $\rho = \varphi
\rho_{\text{eau}} + (1 - \varphi)\rho_{\text{grain}}$, ce qui donne la formule~; $(2710 -
2573)/(2710 - 1000) = 0.080$. (b) $0.51\,e^{-956/4545} = 0.413$. (c) Cimentation~: de la calcite précipitée
dans les pores, que subissent aussi les carbonates en plus de la compaction, retire bien plus de porosité
que le seul enfouissement.
\end{solution}

\begin{solution}{exo:b2:palaeoenvironments:9}
(a) Turbidites de mer profonde (\omterm{def:b2:palaeoenvironments:ichnofacies}{ichnofaciès} de Nereites)~; un fond jurassique du plateau inférieur ou du
talus (Zoophycos)~; un fond au large, entre les bases de vagues (Cruziana)~; un avant-côte (Skolithos)~; un
marais continental. (b) L'eau est devenue de moins en moins profonde vers le haut~: le rivage a avancé vers la
mer. (c) Une baisse du niveau marin, ou un apport de sédiment plus rapide que la subsidence, comme devant
un delta en progression~; une baisse se retrouverait au même moment dans d'autres bassins, alors qu'un delta
serait local et son sédiment viendrait d'un seul côté. (d) Dans une succession concordante, les milieux voisins
se succèdent~; un saut par-dessus les milieux intermédiaires signale une lacune de temps ou une faille.
\end{solution}

\begin{solution}{exo:b2:palaeoenvironments:10}
(a) Le temps d'enfouissement vaut $h/s$ = 100, 10 et \qty{1}{an}. Corps mou : $e^{-100} \approx
\num{4e-44}$, $e^{-10} = \num{4.5e-5}$, $e^{-1} = 0.37$. Coquille ($k = \qty{0.01}{an^{-1}}$)~:
$e^{-1} = 0.37$, $e^{-0.1} = 0.90$, $e^{-0.01} = 0.99$. (b) $s = kh/\ln 100 = 0.1/4.605 =
\qty{0.022}{m.an^{-1}}$, soit \qty{22}{mm} par an. (c) Difficilement, sur des mois~; mais un lit de tempête, une
turbidite ou une chute de cendres enfouit le fond en quelques heures. (d) Le taux de sédimentation constant~:
le temps d'enfouissement est alors bien plus court que $h/s$ avec le taux moyen, et $P = e^{-kT}$, où $T$
est le temps réel d'enfouissement, tend vers un lorsque $T$ est bien inférieur à $1/k$.
\end{solution}

\begin{solution}{exo:b2:palaeoenvironments:11}
(a) Une fine couche $\mathrm dz$ à la profondeur $z$ contient les grains $(1 - \varphi(z))\,\mathrm dz$~;
en surface, elles remplissent $\mathrm dz_0$ avec $(1 - \varphi_0)\,\mathrm dz_0 = (1 -
\varphi(z))\,\mathrm dz$. (b) $\varphi = 0.63\,e^{-816/1960} = 0.415$~; $(1 - 0.415)/0.37 =
1.581$, soit \qty{158.1}{m}. (c) \qty{9.5}{m} de trop~: comme dans l'\cref{exo:b2:palaeoenvironments:7},
la couche restaurée n'est pas uniformément aussi poreuse que sa surface. (d) Le taux de grains,
$h(1-\varphi)/\Delta t$, n'utilise que l'épaisseur et la porosité actuelles, toutes deux mesurées.
\end{solution}

\begin{solution}{exo:b2:palaeoenvironments:12}
(a) $9.5/0.165 = 57.6$ horizons. (b) $180/57.6 = \qty{3.1}{m}$ de calcaire aujourd'hui, $338/57.6 =
\qty{5.9}{m}$ de sédiment frais. (c) Au plus \qty{1.5}{m}~: $180/1.5 = 120$ intervalles, soit 121
échantillons. (d) Des lacunes de dépôt, des lits où les ammonites sont absentes, et des horizons condensés en
quelques centimètres là où le sédiment arrivait lentement.
\end{solution}

\begin{solution}{pb:b2:palaeoenvironments:1}
\textbf{1.} $866 \sin \ang{12} = 866 \times 0.2079 = \qty{180.0}{m}$.
\textbf{2.} $481 \times 0.2079 = \qty{100.0}{m}$. \textbf{3.} Le grès, par
superposition~: les flancs ne plongent que de \ang{12}, les couches ne sont donc
pas renversées, et la plus basse est la plus ancienne. \textbf{4.}
$866 - 766 = \qty{100}{m}$ et $1046 - 866 = \qty{180}{m}$. \textbf{5.} Au moins
\qty{766}{m} à cet endroit~; son sommet d'origine est érodé, si bien que son
épaisseur totale est inconnue. \textbf{6.} 0.70 en surface~;
$0.70\,e^{-766/1408} = 0.406$ à \qty{766}{m}. \textbf{7.} Ce sont des groupes sténohalins~:
l'eau avait la salinité de la pleine mer. \textbf{8.} Un \omterm{def:b2:palaeoenvironments:ichnofacies}{ichnofaciès}
à Cruziana~: le large, entre la base de la vague de beau temps et celle de
la vague de tempête. \textbf{9.} Des tempêtes~: leurs vagues atteignaient
le fond, brisaient les coquilles et les déposaient en lits à base érosive.
\textbf{10.} L'eau de fond était pauvre en oxygène~: peu de fouisseurs,
laminations conservées, pyrite formée. \textbf{11.} L'\omterm{def:b2:palaeoenvironments:ichnofacies}{ichnofaciès} de Zoophycos,
au Mésozoïque sur le plateau inférieur ou le talus, sous la base de la
vague de tempête~; de l'oxygène suffisant par moments pour les fouisseurs.
\textbf{12.} Non~: les ammonites nageaient dans l'eau au-dessus~; leurs
coquilles ont coulé ou dérivé jusqu'au fond après la mort et forment une
\omterm{def:b2:palaeoenvironments:assemblages}{thanatocénose}. \textbf{13.} $184.2 - 174.7 = \qty{9.5}{Ma}$. \textbf{14.}
$180/9.5 = \qty{18.9}{m}$ par million d'années. \textbf{15.}
$(2710 - 2573)/(2710 - 1000) = 0.080$. \textbf{16.} $180 \times (1 - 0.080)/(1 - 0.51) = \qty{338}{m}$.
\textbf{17.} $338/9.5 = \qty{35.6}{m}$ par million d'années.
\textbf{18.} $180 \times 0.92/9.5 = \qty{17.4}{m}$ par million d'années.
\textbf{19.} L'argilite donne 11.3 et \qty{4.4}{m} par million d'années~:
le taux des grains a chuté d'environ quatre fois. Le carbonate du haut-fond
et du plateau a cessé d'arriver à mesure que le fond s'approfondissait,
et seule une lente pluie de boue l'a atteint. \textbf{20.} La mer a avancé~:
des rivières, puis un haut-fond, un fond au large, un fond sous la base
de la vague de tempête, puis la mer ouverte de la craie. La colonne
enregistre une transgression. \textbf{21.} $57 - 41 = \qty{16}{m}$ sur le calcaire,
$86 - 57 = \qty{29}{m}$ sur l'argilite, \qty{45}{m} au total.
\textbf{22.} $(60 - 5) + 338 + 148.6 = \qty{542}{m}$, en
$184.2 - 161.5 = \qty{22.7}{Ma}$~: \qty{23.9}{m} par million d'années. \textbf{23.}
$45/542 = 0.083$~: environ \qty{8}{\percent}, et moins si la courbe des cycles
exagère l'amplitude. Le reste vient de la subsidence du fond. \textbf{24.} La
boue calcaire du calcaire de Varden s'est accumulée à environ
\textbf{\qty[text-series-to-math=true]{36}{m}} par million d'années lors de son dépôt.
\end{solution}
